% !!! 本文件由 work/tools/gen_index.py --phase final 生成，请勿手改 !!!
% 词条按拼音排序（Windows Chinese_China.936 排序规则）；页码取自 main.aux。

\begin{multicols}{2}

\idxhead{符号与数字}
\idxmain{[Boost].UT}{293、294}
\idxmain{4+1 模型}{81、82、86}
\idxsub{场景}{85}
\idxsub{进程视图}{83}
\idxsub{开发视图}{84}
\idxsub{逻辑视图}{82}
\idxsub{物理视图}{85}
\idxhead{A}
\idxmain{AddressSanitizer（ASan）}{358}
\idxmain{AddressSanitizerLeakSanitizer（LSan）}{358}
\idxmain{Aedes}{426}
\idxmain{Alloy}{542}
\idxmain{Amazon Elastic Compute Cloud（EC2）}{531}
\idxmain{Amazon Kinesis}{426}
\idxmain{Amazon MQ}{426}
\idxmain{Amazon Simple Storage Service（S3）}{439}
\idxmain{Ansible}{}
\idxsub{集成到 CI/CD 流水线}{322}
\idxsub{使用 playbook}{322}
\idxsub{使用可复用组件（角色）}{322}
\idxmain{ansible-bender}{519}
\idxmain{Apache ActiveMQ}{425}
\idxmain{Apache Camel}{422}
\idxmain{Apache Kafka}{425、460、463}
\idxsub{系统架构}{461}
\idxmain{Apache Kafka 集群管理器（CMAK）}{462}
\idxmain{Apache Qpid}{425}
\idxmain{Apache ServiceMix}{422}
\idxmain{Apache Synapse}{422}
\idxmain{Apache Thrift}{493、494}
\idxmain{API}{}
\idxsub{管理}{142}
\idxsub{管理工具}{142、143}
\idxsub{网关}{143}
\idxmain{API 调用}{}
\idxsub{通过 CLI 工具使用}{445}
\idxsub{直接使用}{443、445}
\idxmain{API 管理（APIM）}{142}
\idxmain{API 描述语言}{486}
\idxsub{AsyncAPI}{490、492}
\idxsub{OpenAPI}{486、489}
\idxmain{Application Armor（AppArmor）}{508}
\idxmain{Application Verifier（AppVerifier）}{357}
\idxmain{ArgoCD}{592}
\idxmain{AsyncAPI}{490、492}
\idxmain{ATA（高级技术附件，Advanced Technology Attachment）over Ethernet（AoE）}{439}
\idxmain{Auditbeat}{538}
\idxmain{AWS ECS}{531}
\idxmain{AWS Fargate}{531}
\idxmain{Azure Container Registry（ACR）}{513}
\idxmain{Azure Service Bus}{426}
\idxmain{Azure Service Fabric}{531}
\idxmain{安全计算模式（seccomp）}{508}
\idxmain{安全类型}{}
\idxsub{使用}{177}
\idxmain{安全套接字层/传输层安全（SSL/TLS）}{455}
\idxmain{安全外壳（SSH）}{589}
\idxmain{安全意识设计}{338}
\idxsub{GSL}{346}
\idxsub{编码指南}{346}
\idxsub{并发缺陷}{344、346}
\idxsub{防御式编码}{347、350}
\idxsub{接口，设计}{339、342}
\idxsub{漏洞}{351}
\idxsub{自动资源管理，启用}{342、343}
\idxmain{安全增强型 Linux（SELinux）}{508}
\idxhead{B}
\idxmain{Binary JSON（BSON）}{134、468}
\idxmain{Black Duck}{}
\idxsub{参考链接}{356}
\idxmain{BlazingMQ}{426}
\idxmain{Boehm-Demers-Weiser 垃圾回收器（bdwgc）}{342}
\idxmain{Boost.Cobalt}{}
\idxsub{使用示例}{407、408}
\idxmain{Boost.Redis}{}
\idxsub{参考链接}{437}
\idxmain{Boost.Test}{292}
\idxmain{Buildah}{518、519}
\idxmain{buildx}{517}
\idxmain{Business Source License（BSL）}{335}
\idxmain{巴士系数}{112}
\idxmain{包管理}{258}
\idxmain{包含目录}{238}
\idxmain{背压模式}{121}
\idxmain{本地部署模型}{106}
\idxmain{比较并交换（CAS）}{345}
\idxmain{边车模式}{135、136}
\idxsub{实现，使用 Envoy}{136、139}
\idxmain{边车设计模式}{}
\idxsub{用于故障检测}{122}
\idxmain{编写架构文档}{81}
\idxsub{4+1 模型}{81、86}
\idxsub{C4 模型}{86、90}
\idxsub{在敏捷项目中}{91}
\idxmain{编写需求文档}{77}
\idxsub{版本历史，管理}{79}
\idxsub{范围，编写文档}{78}
\idxsub{非功能性需求（NFR），编写文档}{79}
\idxsub{功能性需求，编写文档}{78、79}
\idxsub{上下文，编写文档}{77}
\idxsub{在敏捷项目中}{80}
\idxsub{章节}{80}
\idxmain{编译后的模块接口（CMI）}{281}
\idxmain{编译器}{226}
\idxsub{发现潜在的代码问题}{231、233}
\idxsub{减少构建时间}{228、231}
\idxmain{编译器友好的代码}{}
\idxsub{编写缓存友好的代码}{388、390}
\idxsub{基于真实使用模式进行优化}{388}
\idxsub{设计}{390、392}
\idxsub{生成}{387}
\idxsub{优化整个程序}{388}
\idxmain{标准并行算法}{}
\idxsub{使用}{396}
\idxmain{标准模板库（STL）}{33、149}
\idxmain{表述性状态转移（REST）}{52、473、477}
\idxsub{约束，用于 Web 服务实现}{473}
\idxmain{并发安全的数据结构}{346}
\idxmain{并发视图}{93}
\idxmain{不可变基础设施}{331}
\idxsub{编排，使用 Terraform}{334、335}
\idxsub{实际考量}{331}
\idxsub{实例镜像，使用 Packer 构建}{332、333}
\idxsub{优势}{331}
\idxmain{不可复制类型}{}
\idxsub{实现}{185}
\idxmain{不透明 typedef}{342}
\idxmain{不要重复自己（DRY）}{43}
\idxmain{部署}{563}
\idxmain{部署策略，微服务}{}
\idxsub{每主机单服务}{434}
\idxsub{每主机多服务}{434}
\idxmain{部署代码}{}
\idxsub{构建}{323}
\idxmain{部署即代码}{321}
\idxsub{Ansible 集成，到 CI/CD 流水线}{322}
\idxsub{playbook，在 Ansible 中使用}{322}
\idxsub{管理}{321}
\idxsub{可复用组件（角色），在 Ansible 中使用}{322}
\idxsub{使用}{321}
\idxmain{部署视图}{94、95}
\idxhead{C}
\idxmain{C++}{}
\idxsub{哲学}{31、34}
\idxsub{指定容器接口}{149、153}
\idxmain{C++ Server Pages（CSP）}{348}
\idxmain{C++ 模块}{271、272}
\idxmain{C4 模型}{86、90}
\idxsub{代码图}{89}
\idxsub{容器图}{87、88}
\idxsub{组件图}{88}
\idxmain{Catch2}{291}
\idxsub{使用其进行基准测试}{377、379}
\idxsub{示例}{291}
\idxmain{change root（chroot）}{503}
\idxmain{Checkmarx SAST}{}
\idxsub{参考链接}{355}
\idxmain{Cheerp}{}
\idxsub{参考链接}{359}
\idxmain{CheerpJ}{}
\idxsub{参考链接}{359}
\idxmain{CheerpX}{}
\idxsub{参考链接}{359}
\idxmain{Chiel 法则}{229}
\idxmain{CI/CD}{}
\idxsub{为基础设施验证自动化}{306}
\idxsub{为应用测试自动化}{305、306}
\idxmain{CI/CD 流水线}{}
\idxsub{把 Ansible 集成进去}{322}
\idxmain{CircleCI}{312}
\idxmain{CMake}{}
\idxsub{创建项目}{235}
\idxsub{使用}{234}
\idxsub{用它构建容器镜像}{520、521}
\idxsub{用它配置 Dockerfile}{520}
\idxsub{与模块一起使用}{278、279}
\idxsub{与容器集成}{520}
\idxsub{指定目标}{236、237}
\idxsub{指定输出目录}{237}
\idxmain{CMake 目录变量}{}
\idxsub{区分}{235}
\idxmain{CMake 与高质量代码}{}
\idxsub{安装}{247、250}
\idxsub{导出}{250}
\idxsub{复用}{247}
\idxsub{使用 CPack}{250、254}
\idxmain{CodeSonar}{}
\idxsub{参考链接}{355}
\idxmain{Conan}{313}
\idxsub{与模块一起使用}{279、281}
\idxmain{Conan 包}{264}
\idxsub{测试}{264、265}
\idxsub{创建}{265}
\idxsub{发布}{265、266}
\idxmain{Conan 包管理器}{258}
\idxsub{安装 Conan 依赖}{262}
\idxsub{从 CMake 中使用 Conan 目标}{263}
\idxsub{使用}{258}
\idxsub{指定 Conan 依赖}{260、262}
\idxsub{准备 Conan 配置文件}{259}
\idxmain{Connextra 模板}{71}
\idxmain{const}{}
\idxsub{使用}{177}
\idxmain{Consul 服务网格}{587、588}
\idxmain{coost}{}
\idxsub{使用，示例}{406}
\idxmain{copy-and-swap 惯用法}{}
\idxsub{异常安全，提供}{188、189}
\idxmain{CoroKafka}{400}
\idxmain{CPack}{}
\idxsub{使用}{250、254}
\idxmain{Cppcheck}{}
\idxsub{参考链接}{355}
\idxmain{CppDepend}{95}
\idxmain{CppUMock}{296、301、302}
\idxmain{CppUTest}{291}
\idxmain{CronJobs}{563}
\idxmain{CVE 编号机构（CNA）}{352}
\idxmain{参数化多态}{340}
\idxmain{参数依赖查找（ADL）}{187}
\idxmain{舱壁模式}{124}
\idxmain{操作系统（OS）}{451}
\idxmain{测试金字塔}{285、286}
\idxmain{测试驱动的类设计}{303}
\idxsub{测试，编写}{303}
\idxsub{测试与类设计的冲突}{303}
\idxsub{防御式编程}{304、305}
\idxmain{测试驱动开发（TDD）}{303}
\idxmain{测试驱动自动化}{}
\idxsub{测试，为 CI 编写}{320}
\idxsub{持续测试}{320}
\idxsub{探索}{318}
\idxsub{行为驱动开发}{318、319}
\idxmain{测试替身}{295}
\idxsub{CppUMock 示例}{301、302}
\idxsub{gMock 示例}{296、299}
\idxsub{trompeloeil 示例}{300、301}
\idxsub{编写}{296}
\idxmain{查询语言}{481}
\idxmain{超媒体作为应用状态引擎（HATEOAS）}{473}
\idxmain{超文本传输协议（HTTP）}{338}
\idxmain{超文本应用语言（HAL）}{479}
\idxmain{池资源}{}
\idxsub{使用}{220}
\idxmain{持续部署（CD）}{323}
\idxsub{CD 流水线示例，使用 GitLab 构建}{324、325}
\idxsub{pre-commit 工作流，使用 GitHub Actions 实现}{325、327}
\idxsub{流水线，使用 GitHub Actions 实现}{327、330}
\idxsub{实现}{323、324}
\idxmain{持续测试}{320}
\idxmain{持续集成（CI）}{311}
\idxsub{短发布周期}{311}
\idxsub{流水线，使用 GitLab 实现}{312、315}
\idxsub{实现}{311}
\idxsub{优点}{312}
\idxmain{持续集成/持续部署（CI/CD）}{588}
\idxmain{抽象语法记法一（ASN.1）}{506}
\idxmain{传感器网络 MQTT（MQTT-SN）}{454}
\idxmain{传输层安全（TLS）}{566}
\idxmain{传输控制协议（TCP）}{436}
\idxmain{传输控制协议/网际协议（TCP/IP）}{451}
\idxmain{创建、读取、更新、删除（CRUD）}{62、127、129、474}
\idxmain{垂直扩展}{431}
\idxmain{存储区域网络（SAN）}{439}
\idxmain{存活时间（TTL）}{133、464}
\idxhead{D}
\idxmain{DaemonSet}{563}
\idxmain{Debian 软件包}{258}
\idxmain{DevSecOps}{365}
\idxmain{Distrobox}{359}
\idxmain{Docker Swarm}{527}
\idxmain{Dockerfile}{}
\idxsub{配置，使用 CMake}{520}
\idxsub{容器镜像，使用其构建}{511、512}
\idxmain{Doctest}{292}
\idxmain{Doxygen}{96}
\idxmain{DrogonTest}{293}
\idxmain{代理（broker）}{455}
\idxmain{代理式消息系统}{460}
\idxsub{Apache Kafka}{460、463}
\idxsub{RabbitMQ}{463、466}
\idxmain{代码测试}{}
\idxsub{测试金字塔}{285、287}
\idxsub{对其的需求}{285}
\idxsub{非功能性测试}{287}
\idxsub{根本原因分析}{288}
\idxsub{回归测试}{287}
\idxsub{预防措施}{288、289}
\idxmain{代码加固}{353}
\idxsub{进程隔离}{359、362}
\idxsub{面向安全的内存分配器}{353、354}
\idxsub{沙箱}{362}
\idxsub{自动化检查}{354}
\idxmain{代码库}{}
\idxsub{迁移到模块}{277}
\idxmain{代码评审}{316}
\idxsub{不同方法}{317}
\idxsub{使用拉取请求（合并请求）}{317}
\idxmain{代码图}{89}
\idxmain{单编译单元构建}{230}
\idxmain{单调内存资源}{}
\idxsub{使用}{219}
\idxmain{单态化}{340}
\idxmain{单体仓库（mono-repo）}{429}
\idxmain{单体应用}{}
\idxsub{优缺点}{54}
\idxmain{单一定义规则（ODR）}{271}
\idxmain{单一职责原则（SRP）}{34、51}
\idxmain{单指令多数据（SIMD）}{386}
\idxmain{单子接口}{}
\idxsub{使用}{163、166}
\idxmain{地道 C++}{}
\idxsub{编写}{184}
\idxmain{地址空间布局随机化（ASLR）}{358}
\idxmain{低开销消息系统}{453}
\idxsub{MQTT}{454、457}
\idxsub{ZeroMQ}{458、459}
\idxmain{迪米特法则}{35}
\idxmain{动态多态}{193}
\idxmain{动态分配}{}
\idxsub{减少，使用 SSO/SBO/SOO}{213、215}
\idxmain{动态分析}{}
\idxsub{Application Verifier}{357}
\idxsub{sanitizer}{358}
\idxsub{Valgrind}{357}
\idxsub{模糊测试}{358}
\idxmain{动态链接库（DLL）文件}{237}
\idxmain{动态视图加载（DVL）}{348}
\idxmain{动态应用安全测试（DAST）}{357}
\idxmain{读穿透缓存}{131}
\idxmain{读-改-写（RMW）}{345}
\idxmain{端口}{66}
\idxmain{短字符串优化（SSO）}{213}
\idxmain{断路器模式}{123}
\idxmain{对象}{}
\idxsub{创建}{200}
\idxmain{对象关系映射（ORM）}{261}
\idxmain{多编译器}{}
\idxsub{使用}{227、228}
\idxmain{多仓库（multi-repo）}{429}
\idxmain{多对多通信}{452}
\idxmain{多对一通信}{452}
\idxmain{多态分配器}{}
\idxsub{使用}{216}
\idxmain{多态内存资源（PMR）}{216}
\idxhead{E}
\idxmain{Eclipse Mosquitto MQTT Broker}{425}
\idxmain{ejabberd}{426}
\idxmain{Elastic Container Registry（ECR）}{513}
\idxmain{Elastic Container Service（ECS）}{531}
\idxmain{Elasticsearch}{540}
\idxmain{Elasticsearch、Logstash、Kibana（ELK）}{112}
\idxmain{EMQX}{426}
\idxmain{Enterprise Architect}{96}
\idxmain{Envoy}{}
\idxsub{用于实现边车模式}{136、139}
\idxmain{Envoy Proxy}{}
\idxsub{参考链接}{121}
\idxmain{二进制模块接口（BMI）}{280}
\idxhead{F}
\idxmain{FakeIt}{296}
\idxmain{Fil 的不可思议垃圾收集器（FUGC）}{343}
\idxmain{Filebeat}{538}
\idxmain{find 脚本}{}
\idxsub{编写}{241、243}
\idxsub{使用}{241}
\idxmain{Firecracker}{}
\idxsub{参考链接}{359}
\idxmain{Flawfinder}{}
\idxsub{参考链接}{355}
\idxmain{Fluent Bit}{539}
\idxmain{Fluentd}{538}
\idxmain{FluxCD}{592}
\idxmain{FreeGuard}{}
\idxsub{参考链接}{353}
\idxmain{反馈导向优化（FDO）}{388}
\idxmain{范畴论}{164}
\idxmain{防腐层模式}{126}
\idxmain{防御式编程}{304}
\idxmain{非功能性测试}{287}
\idxmain{非功能性需求（NFR）}{70、71}
\idxsub{约束}{72}
\idxsub{质量属性}{72}
\idxmain{非良构无需诊断（IFNDR）}{274}
\idxmain{非统一内存访问（NUMA）}{94、386}
\idxmain{分布式计算}{}
\idxsub{缺陷，避免}{110、112}
\idxmain{分布式拒绝服务（DDoS）}{482、557}
\idxmain{分布式系统}{}
\idxsub{CAP 定理}{113}
\idxsub{分布式计算缺陷，避免}{110}
\idxsub{服务模型}{106}
\idxsub{特征}{106}
\idxsub{最终一致性}{113}
\idxmain{分层架构}{63、64}
\idxsub{服务于前端的后端（BFF）}{64}
\idxmain{服务}{}
\idxsub{属性}{421}
\idxmain{服务}{563}
\idxmain{服务发现（SD）}{558}
\idxmain{服务级别协议（SLA）}{117}
\idxmain{服务间通信（ISC）}{424、450、451}
\idxmain{服务模型}{106}
\idxsub{本地部署模型}{106}
\idxsub{函数即服务（FaaS）}{109}
\idxsub{基础设施即服务模型（IaaS）}{107}
\idxsub{平台即服务模型（PaaS）}{108}
\idxsub{软件即服务（SaaS）}{108}
\idxmain{服务器端包含（SSI）}{348}
\idxmain{服务器发送事件（SSE）}{492}
\idxmain{服务器消息块（SMB）}{439}
\idxmain{服务网格}{584}
\idxsub{方案}{585}
\idxsub{用于连接服务}{584}
\idxmain{服务网格方案}{}
\idxsub{Consul 服务网格}{587、588}
\idxsub{Istio}{585、586}
\idxsub{Linkerd}{586、587}
\idxmain{服务于前端的后端（BFF）}{64}
\idxmain{服务质量（QoS）}{422、454}
\idxmain{负载均衡器即服务（LBaaS）}{559}
\idxmain{复制模式}{119}
\idxsub{多主复制}{120}
\idxsub{主从复制}{120}
\idxhead{G}
\idxmain{GCC}{227}
\idxmain{geodes}{124}
\idxmain{Geodes 模式}{124}
\idxmain{GitHub Actions}{312}
\idxmain{GitLab}{312}
\idxmain{GitLab CI}{314}
\idxmain{GitLab CI/CD}{312}
\idxmain{GitOps}{588}
\idxsub{优势}{591}
\idxsub{原则}{589}
\idxmain{GitOps 工具}{592}
\idxsub{ArgoCD}{592}
\idxsub{FluxCD}{592}
\idxsub{Jenkins X}{592}
\idxmain{GitOps 优势}{}
\idxsub{改进的安全性}{591}
\idxsub{更高的生产力}{591}
\idxsub{更高的稳定性与可靠性}{591}
\idxsub{开发者体验}{591}
\idxmain{GitOps 原则}{}
\idxsub{集成，与既有组件}{590}
\idxsub{可审计}{590}
\idxsub{配置漂移预防}{590}
\idxsub{声明式描述}{589}
\idxsub{系统状态在 Git 中版本化}{589}
\idxmain{GNU C 库（glibc）}{512}
\idxmain{GNU 调试器（GDB）}{523}
\idxmain{GNU 隐私卫士（GPG）}{589}
\idxmain{Google Artifact Registry（GAR）}{513}
\idxmain{Google Mock（gMock）}{290、296}
\idxmain{Google 的 ClusterFuzz}{}
\idxsub{参考链接}{358}
\idxmain{GoogleTest（GTest）}{290}
\idxmain{Goss}{306}
\idxsub{用于部署验证}{307}
\idxmain{Grafana}{541}
\idxmain{GraphQL}{482、486}
\idxmain{gRPC 网络操作接口（gNOI）}{496}
\idxmain{gRPC 网络管理接口（gNMI）}{496}
\idxmain{gRPC 远程过程调用（gRPC）}{56、496、499}
\idxmain{gRPC-Web}{496}
\idxmain{GTest}{290}
\idxsub{特性}{290}
\idxmain{高级矢量扩展（AVX）}{386}
\idxmain{高级消息队列协议（AMQP）}{424、463}
\idxmain{高频交易（HFT）}{62}
\idxmain{高性能计算（HPC）}{396}
\idxmain{个人信息保护法（APPI）}{339}
\idxmain{个人信息保护法（PIPL）}{339}
\idxmain{个人信息管理（PIM）}{429}
\idxmain{根因分析}{288}
\idxmain{工厂}{}
\idxsub{返回类型，选择}{203}
\idxsub{使用}{200}
\idxmain{工厂方法}{}
\idxsub{使用}{201}
\idxmain{工厂函数}{}
\idxsub{使用}{202}
\idxmain{工厂类}{}
\idxsub{使用}{203、204}
\idxmain{工作节点}{562}
\idxmain{公共对象请求代理体系结构（CORBA）}{468}
\idxmain{功能测试}{287}
\idxmain{功能视图}{92}
\idxmain{功能性需求}{71}
\idxsub{示例}{71}
\idxmain{构建过程}{}
\idxsub{抽象}{234}
\idxmain{构建器}{}
\idxsub{使用}{205、207}
\idxmain{构建时间}{}
\idxsub{减少}{228}
\idxmain{构建应用容器的替代工具}{}
\idxsub{ansible-bender}{519}
\idxsub{Buildah}{518}
\idxmain{故障检测}{122}
\idxsub{边车设计模式，用于}{122}
\idxsub{故障影响，最小化}{123、124}
\idxsub{漏桶计数器模式，用于}{122}
\idxsub{心跳机制模式，用于}{122}
\idxmain{故障转移}{118}
\idxmain{关系数据库服务（RDS）}{531}
\idxmain{关系数据库管理系统（RDBMS）}{114}
\idxmain{管道}{124}
\idxsub{与过滤器模式}{124、125}
\idxmain{管理插件}{465}
\idxmain{惯用法}{184}
\idxmain{过滤器}{124}
\idxhead{H}
\idxmain{hardened\_malloc}{}
\idxsub{参考链接}{353}
\idxmain{HashiCorp 配置语言（HCL）}{528}
\idxmain{HATEOAS}{478、480}
\idxmain{Heartbeat}{538}
\idxmain{Hiredis}{}
\idxsub{参考链接}{437}
\idxmain{HiveMQ}{426}
\idxmain{HTTP}{467}
\idxmain{Hyperlight}{}
\idxsub{参考链接}{359}
\idxmain{函数即服务（FaaS）模型}{109}
\idxmain{黑盒}{427}
\idxmain{后进先出（LIFO）}{401}
\idxmain{后置条件}{}
\idxsub{指定}{159、160}
\idxmain{互斥（mutex）}{343}
\idxmain{互联网小型计算机系统接口（iSCSI）}{439}
\idxmain{环境加固}{363}
\idxsub{DevSecOps}{364、365}
\idxsub{安全影响}{363、364}
\idxsub{策略，链接}{363、364}
\idxsub{地址空间布局随机化}{364}
\idxmain{缓存}{}
\idxsub{更新}{131、133}
\idxmain{缓存模式}{130}
\idxmain{缓存友好的代码}{}
\idxsub{编写}{388、389}
\idxmain{回归测试}{287}
\idxmain{回归测试选择（RTS）}{288}
\idxmain{会话发起协议（SIP）}{426}
\idxmain{混合部署}{107}
\idxhead{I}
\idxmain{IBM Integration Bus}{422}
\idxmain{Istio}{585、586}
\idxhead{J}
\idxmain{Jaeger}{550}
\idxmain{Jakarta Messaging（JMS）API}{425}
\idxmain{Java 服务器页面（JSP）}{348}
\idxmain{Java 虚拟机（JVM）}{359}
\idxmain{JavaScript 对象表示法（JSON）}{471}
\idxsub{数据类型}{471}
\idxmain{JBoss ESB}{422}
\idxmain{Jenkins}{312}
\idxmain{Jenkins X}{312、592}
\idxmain{Jenkins 配置即代码（JCasC）}{321}
\idxmain{Job}{563}
\idxmain{基本可用、软状态、最终一致性（BASE）}{114}
\idxmain{基础设施即代码（IaC）}{447、560}
\idxmain{基础设施即服务（IaaS）模型}{107}
\idxmain{基础设施验证自动化}{}
\idxsub{测试，使用 Goss}{307}
\idxsub{测试，使用 Serverspec}{307}
\idxsub{测试，使用 Testinfra}{307}
\idxsub{用于 CI/CD}{306}
\idxmain{基于 JSON 的 Web 服务}{471、472}
\idxsub{JSON-RPC}{472}
\idxmain{基于 XML 的 Web 服务}{468}
\idxsub{SOAP}{469、470}
\idxsub{UDDI}{470}
\idxsub{XML-RPC}{468}
\idxsub{XML-RPC，与 SOAP 的关系}{468}
\idxmain{基于编排的 saga 模式}{115、117}
\idxmain{基于策略的设计惯用法}{}
\idxsub{使用}{192}
\idxmain{基于队列的负载均衡模式}{121}
\idxmain{基于命令的 saga 模式}{115}
\idxmain{基于模块的架构}{65、66}
\idxmain{基于内核的虚拟机（KVM）}{530}
\idxmain{基于事件的 saga 模式}{115}
\idxmain{基于事件的架构}{60、61}
\idxmain{基于协同（choreography）的 saga 模式}{115}
\idxmain{集成可观测性方案}{552}
\idxmain{集合}{474}
\idxmain{计算}{}
\idxsub{并行编程的平台与工具}{397}
\idxsub{并行化}{392、395}
\idxsub{迁移到编译期}{175、176}
\idxsub{使用 const}{177}
\idxsub{使用 MPI}{396}
\idxsub{使用 OpenMP}{396、397}
\idxsub{使用标准并行算法}{396}
\idxsub{线程与进程}{395}
\idxmain{计算机应急响应小组（CERT）}{352}
\idxmain{加州消费者隐私法案（CCPA）}{339}
\idxmain{加州隐私权法案（CPRA）}{339}
\idxmain{架构权衡分析方法（ATAM）}{74}
\idxmain{架构重要需求（ASR）}{70、73}
\idxsub{识别}{74}
\idxsub{指标}{74}
\idxmain{监控}{543}
\idxmain{简单对象访问协议（SOAP）}{56、469}
\idxsub{头部}{469}
\idxsub{信封}{469}
\idxsub{正文}{469}
\idxmain{简单面向文本的消息协议（STOMP）}{463}
\idxmain{简化包装器与接口生成器（SWIG）}{440}
\idxmain{简明二进制对象表示法（CBOR）}{468}
\idxmain{健康检查}{}
\idxsub{监控}{548}
\idxmain{交互风格}{451}
\idxsub{多对多通信}{452}
\idxsub{多对一通信}{452}
\idxsub{一对多通信}{452}
\idxsub{一对一通信}{452}
\idxmain{绞杀者模式}{126}
\idxmain{接口}{339}
\idxsub{原始指针与智能指针，使用}{154、158}
\idxmain{接口（IFC）文件}{281}
\idxmain{接口定义语言（IDL）}{493}
\idxmain{接口隔离原则}{38、39}
\idxmain{金丝雀发布模式}{140、142}
\idxmain{进程间通信（IPC）}{450、508}
\idxmain{进程视图}{83}
\idxmain{静态多态}{193}
\idxsub{实现}{194、195}
\idxmain{静态分析}{}
\idxsub{SAST 工具，使用}{355}
\idxsub{悬空引用，使用生命周期约束消除}{356}
\idxmain{静态应用安全测试（SAST）}{355}
\idxhead{K}
\idxmain{Kafka Raft（KRaft）}{461}
\idxmain{Kata Containers}{524}
\idxmain{Kibana}{541}
\idxmain{Klocwork}{}
\idxsub{参考链接}{355}
\idxmain{Koenig 查找}{187}
\idxmain{Kubernetes}{525、527}
\idxsub{部署，方法}{562}
\idxsub{采用，考量}{573}
\idxsub{工作节点}{562}
\idxsub{管理}{563、570}
\idxsub{结构}{560、561}
\idxsub{控制平面}{561}
\idxsub{网络}{571}
\idxsub{用于编排云原生工作负载}{559、560}
\idxmain{Kubernetes Operations（kOps）}{562}
\idxmain{Kubernetes 的 Origin 社区发行版（OKD）}{530}
\idxmain{Kubernetes 对象}{}
\idxsub{CronJobs}{563}
\idxsub{DaemonSet}{563}
\idxsub{Deployment}{563}
\idxsub{Job}{563}
\idxsub{Pod}{563}
\idxsub{Service}{563}
\idxsub{容器}{563}
\idxmain{Kubernetes 网络}{571}
\idxsub{Pod 到 Service 通信}{573}
\idxsub{Pod 间通信}{573}
\idxsub{容器间通信}{572}
\idxsub{外部到内部通信}{573}
\idxmain{开闭原则（OCP）}{35、36}
\idxmain{开发视图}{84、94}
\idxmain{开放多处理（OpenMP）}{396}
\idxmain{开放服务网关规范（OSGi）}{422}
\idxmain{开放计算语言（OpenCL）}{397}
\idxmain{开放式 Web 应用程序安全项目（OWASP）}{}
\idxsub{参考链接}{351}
\idxmain{开放系统互连（OSI）模型}{424、451}
\idxmain{开放组架构框架（TOGAF）}{81}
\idxmain{可编程逻辑控制器（PLC）}{453}
\idxmain{可等待对象}{402}
\idxmain{可复制性}{}
\idxsub{管理}{185}
\idxmain{可观测性，在云原生分布式系统中}{574}
\idxsub{应用，使用 OpenTelemetry 插桩}{574、583}
\idxmain{可扩展标记语言（XML）}{468}
\idxmain{可扩展微服务}{}
\idxsub{设计考量}{435}
\idxmain{可扩展消息与在线状态协议（XMPP）}{426、455}
\idxmain{可扩展性瓶颈，应对}{}
\idxsub{存储，外包}{439}
\idxsub{计算，外包}{439、440}
\idxsub{内存管理，外包}{436}
\idxmain{可移动性}{}
\idxsub{管理}{185}
\idxmain{可移植操作系统接口（POSIX）平台}{227}
\idxmain{可移植异构计算接口（HIP）}{397}
\idxmain{客户端缓存}{131}
\idxmain{客户关系管理（CRM）}{108}
\idxmain{控制组（cgroups）}{507、572}
\idxmain{跨站脚本（XSS）}{350}
\idxmain{快速模拟器（QEMU）}{517}
\idxmain{捆绑包}{66}
\idxmain{扩展，微服务}{}
\idxsub{每主机单服务部署}{440、441}
\idxsub{每主机多服务部署}{441}
\idxmain{扩展伯克利包过滤器（eBPF）}{362、560}
\idxmain{垃圾收集暂停（GC 暂停）}{343}
\idxhead{L}
\idxmain{LGTM}{}
\idxsub{参考链接}{355}
\idxmain{libcoro}{}
\idxsub{使用，示例}{403、405}
\idxmain{libmemcached 库}{436}
\idxmain{Linkerd2}{586}
\idxmain{Linux 容器（LXC）}{503}
\idxmain{LLVM 项目（LLDB）}{523}
\idxmain{Logstash}{537}
\idxmain{Loki}{541、542}
\idxmain{蓝绿部署}{140}
\idxmain{类模板实参推导（CTAD）}{154、210}
\idxmain{类型擦除}{}
\idxsub{使用}{196、197}
\idxmain{里氏替换原则（LSP）}{36、38}
\idxmain{联合文件系统（UFS）}{511}
\idxmain{链接器}{363}
\idxmain{链接时优化（LTO）}{388}
\idxmain{两阶段提交（2PC）}{115}
\idxmain{零法则}{186、187}
\idxmain{零停机部署}{140}
\idxsub{金丝雀发布模式}{140、141}
\idxsub{蓝绿部署模式}{140}
\idxmain{领导者选举模式}{119}
\idxmain{领域驱动设计（DDD）}{30、31}
\idxmain{领域特定语言（DSL）}{422}
\idxmain{流式 SIMD 扩展（SSE）}{386}
\idxmain{流式面向文本的消息协议（STOMP）}{424}
\idxmain{六边形架构}{66、67}
\idxmain{漏桶计数器模式}{}
\idxsub{用于故障检测}{122}
\idxmain{路由信息库接口（gRIBI）}{496}
\idxmain{逻辑视图}{82}
\idxhead{M}
\idxmain{Malamute}{458}
\idxmain{manifest-tool}{}
\idxsub{参考链接}{516}
\idxmain{MathWorks 的 Polyspace Bug Finder}{}
\idxsub{参考链接}{355}
\idxmain{Memcached}{436}
\idxmain{MemorySanitizer（MSan）}{358}
\idxmain{Mend Renovate}{}
\idxsub{参考链接}{289}
\idxmain{Metricbeat}{538}
\idxmain{Microsoft Azure Service Bus}{422}
\idxmain{Microsoft BizTalk Server}{422}
\idxmain{Microsoft Visual C++（MSVC）编译器}{169}
\idxmain{Mimir}{542}
\idxmain{MinIO 演练场部署}{444}
\idxmain{Mirantis Container Runtime（MCR）}{562}
\idxmain{Mozilla 公共许可证（MPL）}{335}
\idxmain{MQTT}{455}
\idxmain{美国国防部架构框架（DoDAF）}{}
\idxsub{参考链接}{81}
\idxmain{门控机制}{315}
\idxsub{探索}{315}
\idxmain{面向 NACK 的可靠多播（NORM）}{458}
\idxmain{面向服务的架构（SOA）}{55、420、421、453}
\idxsub{Web 服务}{423}
\idxsub{架构实现}{421}
\idxsub{企业服务总线（ESB）}{421、423}
\idxsub{挑战}{428}
\idxsub{微服务}{427}
\idxsub{消息代理}{424、425}
\idxsub{消息队列}{424}
\idxsub{优势}{56}
\idxsub{优势}{427、428}
\idxsub{云计算}{426}
\idxmain{面向消息的中间件（MOM）}{422}
\idxmain{敏捷原则}{}
\idxsub{用于开发软件架构}{29、30}
\idxmain{名称服务切换（NSS）}{523}
\idxmain{命令查询分离（CQS）}{130、296}
\idxmain{命令查询职责分离（CQRS）}{62、127、129}
\idxmain{命令行界面（CLI）工具}{445}
\idxsub{通过它调用 API}{445}
\idxmain{命名参数惯用法}{207}
\idxmain{命名构造函数惯用法}{201}
\idxmain{模板参数}{}
\idxsub{约束}{177、180}
\idxmain{模糊测试}{358}
\idxmain{模块}{}
\idxsub{创建}{273、276}
\idxsub{使用}{273}
\idxsub{使用，配合 CMake}{278、279}
\idxsub{使用，配合 Conan}{279、281}
\idxsub{现有代码库，迁移到}{277}
\idxmain{模块测试}{285}
\idxmain{模块单元}{273}
\idxmain{模块分区}{273}
\idxmain{模块链接}{273}
\idxmain{模式定义语言（SDL）}{482}
\idxhead{N}
\idxmain{nanobench}{}
\idxsub{基准测试，使用}{379、380}
\idxmain{NATO 架构框架（NAF）}{}
\idxsub{参考链接}{81}
\idxmain{NATS}{425}
\idxmain{niebloid}{}
\idxsub{编写}{190、191}
\idxmain{Nomad}{528、529}
\idxmain{内存}{}
\idxsub{处理}{212}
\idxsub{闪失}{222、223}
\idxmain{内存缓存}{438}
\idxmain{内存区域}{}
\idxsub{使用}{216、217}
\idxmain{内存资源}{}
\idxsub{编写}{220}
\idxmain{内聚}{46、47}
\idxmain{内联命名空间}{}
\idxsub{采用}{160}
\idxmain{内容分发网络（CDN）}{488}
\idxhead{O}
\idxmain{oneAPI 线程构建块（oneTBB）}{397}
\idxmain{Open Virtuozzo（OpenVZ）}{504}
\idxmain{OpenAPI}{486、489}
\idxmain{OpenESB}{422}
\idxmain{OpenShift}{530}
\idxmain{OpenTelemetry}{551}
\idxsub{用于应用插桩}{574、583}
\idxmain{OpenTelemetry 协议（OTLP）}{551、574}
\idxmain{OpenText Fortify}{}
\idxsub{参考链接}{355}
\idxmain{OpenTracing}{550}
\idxmain{Oracle Enterprise Service Bus}{423}
\idxmain{OWASP Dependency-Check}{}
\idxsub{参考链接}{352}
\idxmain{OWASP Top Ten}{315}
\idxmain{耦合}{44、46}
\idxhead{P}
\idxmain{PACELC 定理}{115}
\idxmain{Packer}{332}
\idxsub{实例镜像，使用其构建}{332、333}
\idxmain{Packetbeat}{538}
\idxmain{Parasoft C/C++test}{}
\idxsub{参考链接}{355}
\idxmain{Peach Fuzzer}{}
\idxsub{参考链接}{358}
\idxmain{placement 语法}{217、218}
\idxmain{pmr 分配器}{221}
\idxmain{Pod}{525、563}
\idxmain{Pod 到 Service 通信}{573}
\idxmain{Pod 间通信}{573}
\idxmain{PortSwigger Burp}{}
\idxsub{参考链接}{358}
\idxmain{pre-commit}{325}
\idxmain{Project Diamond}{73}
\idxmain{Prometheus}{542}
\idxmain{Promtail}{542}
\idxmain{Protocol Buffers（Protobuf）}{112、468}
\idxmain{配置管理（CM）}{590}
\idxmain{片段}{273}
\idxmain{平台即服务（PaaS）}{108}
\idxmain{剖析导向反馈（PDF）}{388}
\idxmain{剖析导向优化（PGO）}{388}
\idxhead{Q}
\idxmain{奇异递归模板模式（CRTP）}{193}
\idxmain{企业对企业（B2B）}{423}
\idxmain{企业服务总线（ESB）}{57、420、421}
\idxsub{缺点}{423}
\idxsub{优势}{423}
\idxmain{企业资源规划（ERP）}{74}
\idxmain{前置条件}{}
\idxsub{指定}{159、160}
\idxmain{嵌入式 DRAM（eDRAM）}{390}
\idxmain{缺点，微服务}{}
\idxsub{额外开销}{433}
\idxsub{更难调试}{433}
\idxsub{依赖，对成熟 DevOps 方法}{433}
\idxhead{R}
\idxmain{RabbitMQ}{425、463、466}
\idxsub{RPC，配合使用}{495}
\idxsub{架构}{464、466}
\idxsub{组件}{464}
\idxmain{RAII 守卫}{}
\idxsub{作用域退出动作，自动化}{184}
\idxmain{Red Hat Enterprise Linux CoreOS（RHCOS）}{530}
\idxmain{Red Hat Fuse}{422}
\idxmain{Redis}{436、437}
\idxsub{数据类型}{436}
\idxmain{redis-cpp 库}{}
\idxsub{参考链接}{437}
\idxmain{redis-plus-plus}{}
\idxsub{参考链接}{437}
\idxmain{Requirements for Jira（R4J）}{95}
\idxmain{ReSharper C++}{}
\idxsub{参考链接}{289}
\idxmain{RESTful 服务}{473}
\idxmain{reStructuredText}{96}
\idxmain{绕写缓存}{133}
\idxmain{热备}{118}
\idxmain{日志记录}{535}
\idxsub{配合微服务}{535}
\idxmain{日志记录，在 C++ 中}{}
\idxsub{使用 spdlog}{535、537}
\idxmain{日志聚合}{540}
\idxsub{Elasticsearch}{540}
\idxsub{Loki}{541}
\idxmain{日志可视化}{541}
\idxsub{Grafana}{541、543}
\idxsub{Kibana}{541}
\idxmain{容错系统}{}
\idxsub{背压模式}{121}
\idxsub{复制模式}{119}
\idxsub{构建}{118}
\idxsub{基于队列的负载均衡模式}{121}
\idxsub{冗余模式}{118、119}
\idxmain{容器}{503、504、563}
\idxsub{多架构，定向}{516、517}
\idxsub{集成，使用 CMake}{520}
\idxsub{缺点}{508}
\idxsub{替代运行时}{524}
\idxsub{应用，场景}{507}
\idxsub{优势}{507、508}
\idxsub{与编译型应用}{513、515}
\idxsub{与容器镜像对比}{510、511}
\idxsub{运行时库}{522、523}
\idxsub{运行时实际考量}{522}
\idxmain{容器编排}{524}
\idxsub{托管服务}{531}
\idxsub{自托管方案}{524}
\idxmain{容器的接口}{}
\idxsub{指定，使用 C++}{149、153}
\idxmain{容器即服务（CaaS）}{107}
\idxmain{容器间通信}{572}
\idxmain{容器接口}{}
\idxsub{指定，使用 C++}{149}
\idxmain{容器镜像}{}
\idxsub{分发}{513}
\idxsub{构建}{509、510}
\idxsub{构建，使用 CMake}{520、521}
\idxsub{构建，使用 Dockerfile}{511、512}
\idxsub{命名}{513}
\idxsub{与容器对比}{510、511}
\idxmain{容器类型}{504、505}
\idxmain{容器图}{87、88}
\idxmain{容器网络接口（CNI）}{571}
\idxmain{冗余模式}{118}
\idxsub{共识，达成}{119}
\idxsub{领导者选举}{119}
\idxmain{入侵检测系统（IDS）}{353}
\idxmain{软件定义网络（SDN）}{65}
\idxmain{软件即服务（SaaS）}{312}
\idxmain{软件架构}{26}
\idxsub{范围}{27}
\idxsub{沟通与文化}{27}
\idxsub{基础}{28}
\idxsub{开发，使用敏捷原则}{29、30}
\idxsub{耦合，与内聚对比}{43}
\idxsub{重要性}{28}
\idxmain{软件架构，基础}{}
\idxsub{利益相关者}{29}
\idxsub{上下文}{29}
\idxsub{业务与技术环境}{29}
\idxmain{软件架构，使用敏捷原则}{}
\idxsub{DDD}{30、31}
\idxmain{软件架构，重要性}{}
\idxsub{技术债务}{28}
\idxsub{意外架构}{28}
\idxhead{S}
\idxmain{SaferCPlusPlus}{}
\idxsub{参考链接}{354}
\idxmain{sanitizer}{358}
\idxmain{SAP Process Integration}{423}
\idxmain{Scudo}{}
\idxsub{参考链接}{353}
\idxmain{Serverspec}{306}
\idxsub{用于部署验证}{307}
\idxmain{SMB（服务器消息块）}{332}
\idxmain{Snyk}{}
\idxsub{参考链接}{289}
\idxmain{SOAP 库}{}
\idxsub{Apache Axis}{471}
\idxsub{gSOAP}{471}
\idxmain{SOLID 原则}{34}
\idxsub{单一职责原则（SRP）}{34}
\idxsub{接口隔离原则}{38、39}
\idxsub{开闭原则（OCP）}{35、36}
\idxsub{里氏替换原则（LSP）}{36、38}
\idxsub{依赖倒置}{39、42}
\idxmain{SonarQube}{}
\idxsub{参考链接}{355}
\idxmain{spdlog}{}
\idxsub{用于 C++ 日志记录}{535、537}
\idxmain{Spring Integration}{422}
\idxmain{SQL 注入}{351}
\idxmain{std::optional}{}
\idxsub{函数参数}{161、163}
\idxsub{函数返回值}{162}
\idxsub{类成员}{163}
\idxsub{应用}{161}
\idxmain{Swagger 规范}{486}
\idxmain{三法则与五法则}{186}
\idxmain{三元组}{135}
\idxmain{沙箱 API（SAPI）}{359}
\idxmain{闪失}{}
\idxsub{内存}{222、223}
\idxmain{设计考量，可扩展微服务}{435}
\idxsub{可扩展性瓶颈，应对}{435}
\idxmain{设计模式，微服务}{}
\idxsub{分解模式}{433}
\idxsub{每服务一数据库模式}{434}
\idxmain{身份与访问管理（IAM）}{503}
\idxmain{审计日志}{535}
\idxmain{审计追踪}{535}
\idxmain{生成器表达式}{}
\idxsub{使用}{238、239}
\idxmain{声明式代码}{}
\idxsub{编写}{167、168}
\idxsub{标准范围，应用}{173、175}
\idxsub{特色商品画廊，展示}{168、172}
\idxmain{十二要素应用方法论}{505、506}
\idxmain{时序数据库（TSDB）}{542}
\idxmain{实时应用（RTA）}{343}
\idxmain{实用通用多播（PGM）}{458}
\idxmain{事件}{209、210}
\idxmain{事件驱动架构（EDA）}{466}
\idxmain{事件溯源模式}{61、62、130}
\idxmain{视图选择，用于文档}{}
\idxsub{并发视图}{93}
\idxsub{部署视图}{94、95}
\idxsub{功能视图}{92}
\idxsub{开发视图}{94}
\idxsub{信息视图}{93}
\idxsub{选择}{92}
\idxsub{运行视图}{94、95}
\idxmain{适配器}{66、136}
\idxmain{受约束应用协议（CoAP）}{455}
\idxmain{属性驱动设计（ADD）}{72}
\idxmain{数据报传输层安全（DTLS）}{455}
\idxmain{数据并行 C++（DPC++）}{397}
\idxmain{数据传输对象（DTO）}{128}
\idxmain{数据存储}{}
\idxsub{NoSQL 技术，选择}{134}
\idxsub{设计}{134}
\idxmain{数据分发服务（DDS）}{455}
\idxmain{数据缓存（D-cache）}{388}
\idxmain{数据库缓存}{131}
\idxmain{双向 TLS（mTLS）}{586}
\idxmain{算法}{}
\idxsub{实现}{409、416}
\idxhead{T}
\idxmain{Tekton}{312}
\idxmain{Tempo}{542}
\idxmain{Terraform}{}
\idxsub{不可变基础设施，使用其编排}{334、335}
\idxmain{Terratest}{306}
\idxmain{Testinfra}{306}
\idxsub{用于部署验证}{307}
\idxmain{ThreadSanitizer（TSan）}{358}
\idxmain{Toolbx}{359}
\idxmain{Travis CI}{312}
\idxmain{Trompeloeil}{296、301}
\idxmain{特色商品画廊}{}
\idxsub{展示}{168、172}
\idxmain{提取、转换、加载（ETL）}{460}
\idxmain{替换失败并非错误（SFINAE）}{178、229}
\idxmain{通信平台即服务（CPaaS）}{108}
\idxmain{通用 GPU（GPGPU）}{397}
\idxmain{通用安全服务应用程序接口（GSSAPI）}{458}
\idxmain{通用可扩展性定律（USL）}{394}
\idxmain{通用漏洞与暴露（CVE）}{352}
\idxmain{通用描述、发现与集成（UDDI）}{470}
\idxmain{通用数据保护条例（GDPR）}{339}
\idxmain{通用唯一标识符（UUID）}{445}
\idxmain{通用职责分配软件模式（GRASP）}{184}
\idxmain{同步通信}{451}
\idxmain{统一构建}{230}
\idxmain{统一计算设备架构（CUDA）}{397}
\idxmain{统一建模语言（UML）}{82}
\idxmain{统一日志}{537}
\idxsub{Filebeat}{538}
\idxsub{Fluent Bit}{539}
\idxsub{Fluentd}{538}
\idxsub{Logstash}{537}
\idxsub{Vector}{539、540}
\idxmain{统一资源标识符（URI）}{469}
\idxmain{统一资源定位符（URL）}{474}
\idxmain{头单元}{273}
\idxmain{图查询语言（GQL）}{481}
\idxsub{数据操作，实现}{482}
\idxmain{图形处理器（GPU）}{383、509}
\idxmain{图形用户界面（GUI）}{453}
\idxmain{推导 this 特性}{}
\idxsub{实现}{198、199}
\idxmain{推导指引}{154}
\idxmain{托管服务}{531}
\idxsub{AWS ECS}{531}
\idxsub{AWS Fargate}{531}
\idxsub{Azure Service Fabric}{531}
\idxsub{使用}{441、442}
\idxhead{U}
\idxmain{UndefinedBehaviorSanitizer（UBSan）}{358}
\idxhead{V}
\idxmain{Valgrind}{357}
\idxmain{Vector}{539}
\idxmain{Vector 重映射语言（VRL）}{539}
\idxmain{Veracode Static Analysis}{}
\idxsub{参考链接}{356}
\idxmain{VerneMQ}{426}
\idxmain{Visual Paradigm}{96}
\idxhead{W}
\idxmain{Web API 设计方法}{}
\idxsub{GraphQL}{481、486}
\idxsub{REST}{473}
\idxmain{Web API 载荷格式}{467}
\idxsub{SOAP 库}{471}
\idxsub{基于 JSON 的 Web 服务}{471、472}
\idxsub{基于 XML 的 Web 服务}{468}
\idxmain{Web 服务}{423}
\idxsub{API 描述语言}{486}
\idxsub{Web API 设计方法}{473}
\idxsub{Web API 载荷格式}{468}
\idxsub{Web API 载荷格式}{467}
\idxsub{工作，工具}{467}
\idxsub{使用}{467}
\idxsub{优缺点}{423}
\idxmain{Web 服务描述语言（WSDL）}{470}
\idxmain{Web 服务器缓存}{131}
\idxmain{WebAssembly（Wasm）}{109、359}
\idxmain{WebVM}{}
\idxsub{参考链接}{359}
\idxmain{Winlogbeat}{538}
\idxmain{外部到内部通信}{573}
\idxmain{外部模块}{}
\idxsub{find 脚本，编写}{241、243}
\idxsub{find 脚本，使用}{241}
\idxsub{测试，添加}{244、247}
\idxsub{使用}{239}
\idxsub{依赖，获取}{239、241}
\idxmain{外设接口控制器（PIC）}{453}
\idxmain{外围组件互连（PCI）}{421}
\idxmain{万维网联盟（W3C）}{423、469}
\idxmain{网际协议（IP）}{427}
\idxmain{网络地址转换（NAT）}{525}
\idxmain{网络附加存储（NAS）}{439}
\idxmain{网络块设备（NBD）}{439}
\idxmain{网络文件系统（NFS）}{332、439}
\idxmain{微服务}{56、427、429、505、507}
\idxsub{部署策略}{434}
\idxsub{扩展}{59、440、441}
\idxsub{迁移到}{59}
\idxsub{缺点}{58、432、433}
\idxsub{设计模式}{433}
\idxsub{特征}{57、58}
\idxsub{用于日志记录}{535}
\idxsub{优势}{58、429、432}
\idxmain{微服务的好处}{}
\idxsub{分布式开发}{432}
\idxsub{可扩展性}{431}
\idxsub{灵活性}{432}
\idxsub{模块化}{429、431}
\idxsub{与遗留系统集成}{432}
\idxmain{微基准测试}{}
\idxsub{Catch2，使用}{377、379}
\idxsub{Google Benchmark，设置}{370}
\idxsub{编写}{371、373}
\idxsub{复用，通过参数捕获}{374}
\idxsub{使用}{370}
\idxsub{使用 nanobench}{379、380}
\idxsub{输入规模，通过参数变化}{374、376}
\idxsub{选择，要优化什么}{381}
\idxsub{用于创建性能测试}{381}
\idxmain{位置无关代码（PIC）}{358}
\idxmain{位置无关可执行文件（PIE）}{358}
\idxmain{文档对象模型（DOM）}{360}
\idxmain{文档生成}{95}
\idxsub{API 文档，从代码生成}{96、101}
\idxsub{图表，从代码生成}{95}
\idxsub{需求生成}{95}
\idxmain{无状态方法}{}
\idxsub{与有状态方法对比}{49、51}
\idxmain{无状态服务}{52、53}
\idxmain{物理视图}{85}
\idxmain{物联网（IoT）}{420、507、559}
\idxhead{X}
\idxmain{XML 外部实体（XXE）}{351}
\idxmain{XML-RPC}{468}
\idxmain{系统}{}
\idxsub{集成}{124}
\idxsub{可用性，计算}{117、118}
\idxsub{使其容错}{117}
\idxmain{系统部署}{135}
\idxsub{边车模式}{135}
\idxsub{外部配置存储模式}{142}
\idxmain{系统集成模式}{}
\idxsub{防腐层模式}{126}
\idxsub{管道与过滤器模式}{124、125}
\idxsub{绞杀者模式}{126}
\idxsub{迁移，从遗留模式}{126}
\idxsub{消费者模式，竞争}{125、126}
\idxmain{系统建模语言（SysML）}{}
\idxsub{参考链接}{81}
\idxmain{现场可编程门阵列（FPGA）}{528}
\idxmain{项目管理知识体系（PMBOK）}{73}
\idxmain{消费者模式}{}
\idxsub{竞争}{125、126}
\idxmain{消息传递接口（MPI）}{396}
\idxmain{消息代理}{424、425}
\idxmain{消息队列}{424}
\idxmain{消息队列遥测传输（MQTT）}{424}
\idxmain{消息系统}{453}
\idxsub{代理式消息系统}{460}
\idxsub{低开销消息系统}{453}
\idxmain{小对象优化（SOO）}{213}
\idxmain{小缓冲区优化（SBO）}{213}
\idxmain{协程}{}
\idxsub{Boost.Cobalt，示例}{407、408}
\idxsub{coost，示例}{406、407}
\idxsub{libcoro 实用工具}{401、402}
\idxsub{libcoro，示例}{403、405}
\idxsub{使用}{398、401}
\idxsub{与 awaitable}{402}
\idxmain{写穿透缓存}{132}
\idxmain{写回缓存}{132}
\idxmain{写时复制（COW）}{216}
\idxmain{心跳机制模式}{}
\idxsub{用于故障检测}{122}
\idxmain{新兴架构}{26}
\idxmain{信息视图}{93}
\idxmain{行为驱动开发（BDD）}{291、318、319}
\idxmain{性能测量}{369}
\idxsub{工具，利用}{369}
\idxsub{微基准测试，使用}{370}
\idxsub{性能分析}{382}
\idxsub{追踪}{387}
\idxsub{准备}{369}
\idxmain{性能测试}{}
\idxsub{创建，使用基准测试}{381}
\idxmain{性能分析}{382}
\idxmain{性能分析器}{}
\idxsub{环境，准备}{386}
\idxsub{结果，处理}{386}
\idxsub{结果，分析}{386}
\idxsub{类型，选择}{382、385}
\idxmain{虚拟机（VM）}{331、353、555}
\idxmain{虚拟机监视器（VMM）}{359}
\idxmain{虚拟私有云（VPC）}{445}
\idxmain{虚拟专用服务器（VPS）}{442}
\idxmain{需求技术}{75}
\idxsub{利益相关者}{76}
\idxsub{利益相关者需求}{76、77}
\idxsub{上下文}{75}
\idxsub{现有文档}{76}
\idxmain{需求语法简易方法（EARS）}{78}
\idxhead{Y}
\idxmain{延迟加载}{132}
\idxmain{演进式架构}{29}
\idxmain{业务流程管理（BPM）}{60}
\idxmain{业务流程执行语言（BPEL）}{60}
\idxmain{一对多通信}{452}
\idxmain{一对一通信}{452}
\idxmain{依赖}{}
\idxsub{获取}{239、241}
\idxmain{依赖安全}{352}
\idxsub{CVE 列表，使用}{352}
\idxsub{内部}{352}
\idxsub{外部}{352}
\idxsub{自动化依赖升级管理}{353}
\idxsub{自动扫描器，使用}{352}
\idxmain{依赖倒置}{39、42}
\idxmain{以编译器为中心的工具}{}
\idxsub{使用}{233}
\idxmain{异步通信}{451}
\idxmain{异常安全}{}
\idxsub{提供，使用 copy-and-swap 惯用法}{188、189}
\idxmain{异常跟踪}{540}
\idxmain{异常值检测}{122}
\idxmain{意外架构}{26、28}
\idxmain{引导文件系统（bootfs）}{511}
\idxmain{隐藏友元}{}
\idxsub{编写}{187、188}
\idxmain{应用测试}{}
\idxsub{为 CI/CD 自动化}{305、306}
\idxmain{应用缓存}{131}
\idxmain{应用容器}{}
\idxsub{使用替代工具构建}{518、519}
\idxmain{应用指标}{}
\idxsub{实现}{544、547}
\idxmain{硬件辅助 AddressSanitizer（HWASan）}{358}
\idxmain{用户数据报协议（UDP）}{436}
\idxmain{优化列式存储（ORC）}{468}
\idxmain{有意架构}{28}
\idxmain{有状态方法}{}
\idxsub{与无状态方法对比}{49、51}
\idxmain{有状态服务}{52}
\idxmain{语义属性}{578}
\idxmain{预编译模块（PCM）}{281}
\idxmain{预编译头（PCH）文件}{228}
\idxmain{预刷新缓存}{133}
\idxmain{原型}{}
\idxsub{使用}{207}
\idxmain{原子性、一致性、隔离性、持久性（ACID）}{114}
\idxmain{远程代码执行（RCE）}{348}
\idxmain{远程过程调用（RPC）}{428、468}
\idxsub{Apache Thrift}{493、495}
\idxsub{gRPC}{496、499}
\idxsub{配合 RabbitMQ}{495}
\idxmain{远程字典服务器（Redis）}{425}
\idxmain{约束}{72}
\idxsub{示例}{72}
\idxmain{云 API}{}
\idxsub{访问}{442、443}
\idxsub{使用第三方工具}{446、447}
\idxmain{云 CLI}{}
\idxsub{使用}{446}
\idxmain{云 SDK}{442}
\idxmain{云计算}{426、442}
\idxmain{云提供商}{}
\idxsub{使用}{441、442}
\idxmain{云提供商 API}{}
\idxsub{使用}{443}
\idxmain{云原生}{555}
\idxsub{CNCF 示例}{555}
\idxsub{反向代理和正向代理}{557}
\idxsub{服务发现}{558、559}
\idxsub{负载均衡}{556}
\idxsub{云作为操作系统}{555}
\idxmain{云原生分布式系统}{}
\idxsub{可观测性}{574}
\idxmain{云原生工具}{447}
\idxmain{云原生计算基金会（CNCF）}{525、550、555}
\idxmain{云原生应用保护平台（CNAPP）}{566}
\idxmain{运行时测试框架}{290}
\idxsub{[Boost].UT 示例}{293、294}
\idxsub{Boost.Test 示例}{292}
\idxsub{Catch2 示例}{291}
\idxsub{CppUTest 示例}{291}
\idxsub{Doctest 示例}{292}
\idxsub{DrogonTest 示例}{293}
\idxsub{GTest 示例}{290}
\idxsub{编译期测试技术}{294、295}
\idxmain{运行时库（RTL）}{201}
\idxmain{运行视图}{94、95}
\idxhead{Z}
\idxmain{Zed Attack Proxy（ZAP）}{}
\idxsub{参考链接}{358}
\idxmain{ZeroMQ}{458、460}
\idxmain{ZeroVM}{}
\idxsub{参考链接}{359}
\idxmain{Zipkin}{550}
\idxmain{在线 OpenAPI 编辑器}{}
\idxsub{参考链接}{487}
\idxmain{张量处理单元（TPU）}{530}
\idxmain{长期支持（LTS）}{511}
\idxmain{支付卡行业数据安全标准（PCI DSS）}{586}
\idxmain{指定初始化器}{171}
\idxmain{指令缓存（I-cache）}{388}
\idxmain{指南支持库（GSL）}{343}
\idxmain{质量属性}{72}
\idxsub{示例}{72}
\idxmain{中间表示}{357}
\idxmain{重试模式}{123}
\idxmain{主动-被动（主从）故障转移}{118}
\idxmain{状态}{208}
\idxsub{跟踪}{209、212}
\idxmain{追踪}{549}
\idxsub{Jaeger}{550}
\idxsub{OpenTracing}{550}
\idxmain{资源}{342}
\idxmain{资源获取即初始化（RAII）}{147、148、343}
\idxmain{自动化检查}{}
\idxsub{编译器警告}{354、355}
\idxsub{动态分析}{357}
\idxsub{静态分析}{355}
\idxmain{自托管方案}{524}
\idxsub{Docker Swarm}{527}
\idxsub{Kubernetes}{525、527}
\idxsub{Nomad}{528、530}
\idxsub{OpenShift}{530}
\idxmain{纵深防御原则}{111、288}
\idxmain{组件测试}{285}
\idxmain{组件对象模型（COM）}{468}
\idxmain{组件图}{88}
\idxmain{最少知识原则}{35}
\idxmain{最终用户许可协议（EULA）}{509}
\idxmain{作用域退出动作}{}
\idxsub{自动化，使用 RAII 守卫}{184}

\end{multicols}
